\section*{Uniform shared-data curvature away from the balanced atom}

All logarithms are base two.  Split the cube into $k$ disjoint $r$-bit
blocks and let $f_{k,r}$ be $-1$ exactly when at least one block equals a
fixed point.  Put $\epsilon=2^{-r}$ and expose one physical bit $s$ of
one block.  Its two children share every other physical data coordinate.
Their disagreement probability is
\[
 \delta=\operatorname{Inf}_s(f_{k,r})
 ={2\epsilon(1-\epsilon)^k\over1-\epsilon}.           \tag{1}
\]
For the complete signed GU increment write
\[
 Q_s=G(f_{k,r})-
 {G((f_{k,r})_{s,+})+G((f_{k,r})_{s,-})\over2}
 =\Delta_sU+\sum_{i\ne s}\Delta_s\Delta_iU.          \tag{2}
\]

\paragraph{Theorem.}
Let $K$ be a compact subset of
$(0,\infty)\setminus\{\ln2\}$.  Uniformly for $\lambda\in K$ and
$k_r=\lfloor\lambda2^r\rfloor$, as $r\to\infty$,
\[
 {\Delta_sU\over\delta}
 =1+2r2^{-r}+O_K(2^{-r}),                              \tag{3}
\]
\[
 {Q_s\over\delta}
 =1-2r^2 2^{-r}+O_K(r2^{-r}).                          \tag{4}
\]
Consequently there exists $r_0(K)$ such that for every
$r\ge r_0(K)$ and every $\lambda\in K$,
\[
 \boxed{\Delta_sU>\operatorname{Inf}_s(f_{k_r,r}),
 \qquad Q_s<\operatorname{Inf}_s(f_{k_r,r}).}          \tag{5}
\]
Thus the selector-only charge fails uniformly as the nonuniform branch
gap $\delta\asymp_K2^{-r}$ vanishes, while the complete signed shared-data
increment is paid with coefficient one by the same, nonreused physical
influence.

\paragraph{Exact formulas.}
For a heterogeneous vector of block sizes $\mathbf r$, set
\[
 P=\prod_j(1-\epsilon_j),\quad
 L_j=-\log_2(1-\epsilon_j),\quad
 \epsilon_j=2^{-r_j}.
\]
The exact upper-set calculation is
\[
 U(\mathbf r)=c((2P-1)^2)
 +4P\sum_j\epsilon_j\left(r_j-{r_j\over2(1-\epsilon_j)}\right)
 +4P\sum_jL_j((1-\epsilon_j)-P),                       \tag{6}
\]
where $c(x)=x\log_2(1/x)$.  Restricting $s$ either deletes its block or
replaces it by an $(r-1)$-block.  Therefore (2), including both unequal
children and every normalization term, is obtained by applying (6) first
to $U$ and then to every physical restriction defining $G$.

\paragraph{Selector surcharge.}
Write $P=(1-\epsilon)^k$.  The child products are
\[
 P_-={P\over1-\epsilon},\qquad
 P_+={P(1-2\epsilon)\over1-\epsilon},\qquad
 {P_-+P_+\over2}=P.                                    \tag{7}
\]
For $\lambda\in K$, $P\to e^{-\lambda}$ in a compact subset of
$(0,1)$ disjoint from $1/2$.  Hence $c((2P-1)^2)$ has uniformly bounded
derivatives.  Substitution of (7) in (6), together with
$L=\epsilon/\ln2+O(\epsilon^2)$, gives
\[
 \Delta_sU=\delta\{1+2r\epsilon+O_K(\epsilon)\},       \tag{8}
\]
which is (3).  Block symmetry gives the equivalent exact relation
$G=kr\Delta_sU$ and $I=kr\delta$.

\paragraph{Complete curvature refund.}
Let $E=G-I$.  Influence recursion and commutation of restrictions give
the exact child-residual identity
\[
 Q_s-\delta=E(f_{k,r})-
 {E((f_{k,r})_{s,+})+E((f_{k,r})_{s,-})\over2}.        \tag{9}
\]
Taylor expansion of the complete expression (6), uniform on $K$, yields
for the root and its two one-bit children
\[
 E(\mathbf r)=4P\sum_j r_j^2\epsilon_j^2
                 +\mathcal E(\mathbf r),              \tag{10}
\]
where changing the exposed $r$-block to an $(r-1)$-block or deleting it
changes $\mathcal E$ by $O_K(r\epsilon^2)$.  The leading functional in
(10) has the exact central change
\[
 E(f_{k,r})-
 {E((f_{k,r})_{s,+})+E((f_{k,r})_{s,-})\over2}
 =-4Pr^2\epsilon^2+O_K(Pr\epsilon^2).                 \tag{11}
\]
Dividing (11) by (1) proves (4).

\paragraph{Paid boundary and why this is not closure.}
The root residual itself is positive:
\[
 E(f_{k_r,r})
 =4\lambda e^{-\lambda}r^2 2^{-r}
   +O_K(r2^{-r}).                                      \tag{12}
\]
Rearranging (9) gives the complete stopping ledger
\[
 {E((f_{k,r})_{s,+})+E((f_{k,r})_{s,-})\over2}
 =E(f_{k,r})+(\delta-Q_s).                             \tag{13}
\]
The nonreused reserve $\delta-Q_s$ is genuine and nonnegative by (4),
but (13) transports rather than erases the positive residual: the average
child debt is larger by exactly that reserve.  This is a paid boundary,
not an arbitrary-leaf assumption.  It proves that favorable complete
local drift alone cannot establish $G\le I$.

The excluded intensity $\lambda=\ln2$ is structural.  There
$P\to1/2$, the constant Fourier mass $(2P-1)^2$ vanishes, and its
$x\log(1/x)$ cusp changes the asymptotic scale.  The exact formulas remain
valid, but the bounded-derivative proof of (3)--(4) does not.  No claim at
that exceptional intensity is made here.

\paragraph{GU scope.}
Equations (4), (9), and (13) give a coefficient-one payment for one
complete shared-data restriction step, uniformly over every compact
fixed-intensity range away from the balanced atom.  They do not provide a
universal pair $(a,b)$: a universal proof still needs a potential whose
terminal value pays the transported $E$ debt for arbitrary Boolean
functions without charging any physical influence twice.
